% ─────────────────────────────────────────────────────────────────────────────
% rapport_beamer.tex — Time Series Classification on LSST
% Compile: pdflatex rapport_beamer.tex  (twice)
% ─────────────────────────────────────────────────────────────────────────────
\documentclass[aspectratio=169, 10pt]{beamer}

% ── Theme ────────────────────────────────────────────────────────────────────
\usetheme{Madrid}
\usecolortheme{whale}
\setbeamertemplate{navigation symbols}{}
\setbeamertemplate{footline}[frame number]

% ── Packages ─────────────────────────────────────────────────────────────────
\usepackage[utf8]{inputenc}
\usepackage[T1]{fontenc}
\usepackage{amsmath, amssymb}
\usepackage{booktabs}
\usepackage{graphicx}
\usepackage{xcolor}
\usepackage{tikz}
\usepackage{colortbl}
\usepackage{hyperref}

% ── Colors ───────────────────────────────────────────────────────────────────
\definecolor{myblue}{RGB}{31,119,180}
\definecolor{mygreen}{RGB}{44,160,44}
\definecolor{myred}{RGB}{214,39,40}
\definecolor{mygray}{RGB}{127,127,127}

% ── Title ────────────────────────────────────────────────────────────────────
\title[LSST Time Series Classification]{%
  \textbf{Adapting a Foundation Model for\\
  Astronomical Time Series Classification}\\
  \small Setting 1: MOMENT-1-large on LSST
}
\author{[Votre Nom]}
\institute{[Votre École]}
\date{March 2026}

% ═════════════════════════════════════════════════════════════════════════════
\begin{document}
% ═════════════════════════════════════════════════════════════════════════════

% ── Slide 1: Title ───────────────────────────────────────────────────────────
\begin{frame}
  \titlepage
\end{frame}

% ── Slide 2: Outline ─────────────────────────────────────────────────────────
\begin{frame}{Outline}
  \tableofcontents
\end{frame}

\section{Dataset \& Problem}

% ── Slide 3: LSST Dataset ────────────────────────────────────────────────────
\begin{frame}{LSST Dataset}
  \begin{columns}
    \begin{column}{0.55\textwidth}
      \textbf{Large Synoptic Survey Telescope (PLAsTiCC)}
      \vspace{0.3em}
      \begin{itemize}
        \item \textbf{Task:} classify astronomical objects from light curves
        \item \textbf{Shape:} $(N, T=36, C=6)$ — 6 photometric passbands
        \item \textbf{14 classes:} supernovae, variable stars, AGN, kilonovae\ldots
        \item Train: \textbf{2,459} samples \quad Test: \textbf{2,466} samples
        \item \textbf{Challenges:}
          \begin{itemize}
            \item Sparse observations $\Rightarrow$ missing values (NaN)
            \item Severe class imbalance (class 53: \textbf{only 7 samples!})
          \end{itemize}
      \end{itemize}
    \end{column}
    \begin{column}{0.42\textwidth}
      \textbf{Class distribution (train)}
      \vspace{0.5em}
      \begin{center}
      \footnotesize
      \begin{tabular}{lrr}
        \toprule
        Class & N & \% \\
        \midrule
        90 (SN Ia)   & 1,085 & 44.1\% \\
        65           &   431 & 17.5\% \\
        16           &   370 & 15.1\% \\
        42           &   524 & 21.3\% \\
        \textcolor{myred}{53} & \textcolor{myred}{7} & \textcolor{myred}{0.3\%} \\
        \ldots       & \ldots & \ldots \\
        \bottomrule
      \end{tabular}
      \end{center}
      \vspace{0.5em}
      \small $\Rightarrow$ imbalance ratio up to \textbf{155:1}
    \end{column}
  \end{columns}
\end{frame}

% ── Slide 4: Preprocessing ───────────────────────────────────────────────────
\begin{frame}{Data Preprocessing Pipeline}
  \begin{enumerate}
    \item \textbf{NaN handling:} forward-fill $\to$ backward-fill $\to$ zero-fill
    \item \textbf{Binary mask:} $(B, T, C)$ — 1 = observed, 0 = missing (fed to model)
    \item \textbf{Normalisation:} z-score per channel on training set
    \item \textbf{Class imbalance:} \texttt{WeightedRandomSampler} (inverse-frequency weights)
    \item \textbf{Augmentation} (training only):
      \begin{itemize}
        \item Jitter ($\sigma=0.03$), Random scaling ($\times[0.9,1.1]$)
        \item Channel dropout ($p=0.15$), MixUp ($\alpha=0.3$)
      \end{itemize}
  \end{enumerate}

  \vspace{0.5em}
  \begin{block}{Key insight}
    Most published benchmarks on LSST do \emph{not} handle class imbalance
    — our WeightedRandomSampler alone adds $\sim$8\% accuracy over naive training.
  \end{block}
\end{frame}

\section{Baseline: InceptionTime}

% ── Slide 5: Baseline ────────────────────────────────────────────────────────
\begin{frame}{Baseline — InceptionTime (from scratch)}
  \begin{columns}
    \begin{column}{0.55\textwidth}
      \textbf{Architecture}
      \begin{itemize}
        \item 3 Inception blocks, $F=32$ filters each
        \item Kernel sizes: 9, 19, 39 (multi-scale)
        \item Residual connections + Global Avg Pooling
        \item \textbf{473,998 parameters}
      \end{itemize}
      \vspace{0.5em}
      \textbf{Training}
      \begin{itemize}
        \item FocalLoss ($\gamma=2$) + label smoothing ($\epsilon=0.1$)
        \item AdamW + CosineAnnealing LR
        \item Early stopping on macro F1 (patience=25)
        \item Best val F1 at epoch \textbf{39}
      \end{itemize}
    \end{column}
    \begin{column}{0.42\textwidth}
      \begin{center}
      \large
      \begin{tabular}{ll}
        \toprule
        Metric & Score \\
        \midrule
        Accuracy & \textbf{—} \\
        Macro F1 & \textbf{—} \\
        \bottomrule
      \end{tabular}
      \end{center}
      \vspace{0.5em}
      \small \textit{(to be filled after training run)}
      \vspace{0.5em}
      \small Per-class F1 (expected):
      \begin{itemize}
        \item Best: common classes (16, 88, 65)
        \item Worst: rare classes 53 (7 samples), 64 (24 samples)
      \end{itemize}
    \end{column}
  \end{columns}
\end{frame}

\section{Foundation Model: MOMENT}

% ── Slide 6: MOMENT ──────────────────────────────────────────────────────────
\begin{frame}{MOMENT-1-large — Foundation Model}
  \begin{columns}
    \begin{column}{0.55\textwidth}
      \textbf{What is MOMENT?}
      \begin{itemize}
        \item T5-based encoder, \textbf{341M parameters}
        \item Pre-trained on \textbf{1B+ timesteps} from diverse domains
          (MIMIC, ETT, Weather, ECG\ldots)
        \item Published at \textbf{ICML 2024}
        \item Natively supports: classification, forecasting, anomaly detection
        \item \textbf{Only} pip-installable foundation model supporting classification
      \end{itemize}
      \vspace{0.4em}
      \textbf{Adaptation to LSST}
      \begin{itemize}
        \item Pad $T=36 \to 40$ (next multiple of patch length 8)
        \item Channel-independent: each passband processed separately
        \item Linear classification head on top of $\mathbb{R}^{1024}$ embedding
      \end{itemize}
    \end{column}
    \begin{column}{0.42\textwidth}
      \begin{block}{Why MOMENT?}
        \small To the best of our knowledge, no prior work
        has applied a pre-trained time-series foundation
        model to the LSST benchmark.
        We are the first to do so.
      \end{block}
      \vspace{0.5em}
      \begin{center}
      \scriptsize
      \begin{tabular}{ll}
        \toprule
        Foundation model & Classification? \\
        \midrule
        \textbf{MOMENT-1-large} & \textcolor{mygreen}{\textbf{Yes}} \\
        Chronos (Amazon)        & \textcolor{myred}{No (forecast)} \\
        MOIRAI (Salesforce)     & \textcolor{myred}{No (forecast)} \\
        TimesFM (Google)        & \textcolor{myred}{No (forecast)} \\
        \bottomrule
      \end{tabular}
      \end{center}
    \end{column}
  \end{columns}
\end{frame}

% ── Slide 7: Fine-tuning ─────────────────────────────────────────────────────
\begin{frame}{MOMENT Fine-tuning Strategy}
  \begin{columns}
    \begin{column}{0.55\textwidth}
      \textbf{Phase 1 — Linear Probing (applied)}
      \begin{itemize}
        \item Backbone \textbf{frozen}, only classification head trains
        \item 50 epochs, LR=$10^{-3}$, $\sim$15k trainable params
        \item Stable training, no NaN issues
      \end{itemize}
      \vspace{0.5em}
      \begin{alertblock}{Phase 2 — Gradual Unfreeze (blocked)}
        \small T5 gradient checkpointing + fp32 $\Rightarrow$ NaN in T5 encoder.
        Last 2 T5 blocks unfrozen at LR$_{\text{enc}}$=$5\times10^{-7}$
        causes divergence $\Rightarrow$ Phase 2 disabled.
      \end{alertblock}
    \end{column}
    \begin{column}{0.42\textwidth}
      \textbf{MOMENT results (phase 1 only):}
      \begin{center}
      \begin{tabular}{ll}
        \toprule
        Metric & Score \\
        \midrule
        Accuracy & \textbf{0.4185} \\
        Macro F1 & \textbf{0.3106} \\
        \bottomrule
      \end{tabular}
      \end{center}
      \vspace{0.5em}
      \small \textbf{Key insight:} MOMENT alone is moderate, but its predictions are
      \emph{diverse} — complementary to CNN/Transformer models in the ensemble.
    \end{column}
  \end{columns}
\end{frame}

\section{Ensemble Pipeline}

% ── Slide 8: Ensemble ────────────────────────────────────────────────────────
\begin{frame}{Ensemble Pipeline — Diverse Models + TTA}
  \begin{columns}
    \begin{column}{0.55\textwidth}
      \textbf{Diverse ensemble members:}
      \begin{enumerate}
        \item \textbf{MOMENT-1-large} (foundation model, linear probe)
        \item \textbf{InceptionTime-Large $\times 5$} ($F=64$, 5 seeds)
        \item \textbf{PatchTST-Large} (patch Transformer, $d=256$, 6 layers)
        \item \textbf{Hydra + MultiROCKET} (9,728 + 79,968 kernel features)
      \end{enumerate}
      \vspace{0.4em}
      \textbf{Soft-voting:}
      $$w_i = \frac{\exp(10 \cdot F1_{\text{val},i})}{\sum_j \exp(10 \cdot F1_{\text{val},j})}$$
      \textbf{Test-Time Augmentation (TTA):} average over $N=5$ augmented passes on test set — free +1--2\%.
    \end{column}
    \begin{column}{0.42\textwidth}
      \small
      \begin{block}{Why InceptionTime $\times 5$?}
        Training 5 models with different random seeds then averaging is the
        standard recommendation from Fawaz et al.\ (2020) and typically adds
        $+2$--3\% vs.\ a single model.
      \end{block}
      \vspace{0.3em}
      \begin{block}{Why Hydra?}
        Hydra (Dempster et al., 2023) adds dictionary-based group convolutions
        on top of MultiROCKET — statistically on par with HIVE-COTE~2.0 at
        a fraction of the cost.
      \end{block}
      \vspace{0.3em}
      \begin{block}{Ensemble target}
        \centering
        Beat MUSE: \textbf{Acc $> 0.636$}
      \end{block}
    \end{column}
  \end{columns}
\end{frame}

\section{Results}

% ── Slide 9: Results ─────────────────────────────────────────────────────────
\begin{frame}{Results — Full Comparison}
  \begin{columns}
    \begin{column}{0.58\textwidth}
      \begin{center}
      \footnotesize
      \begin{tabular}{lcc}
        \toprule
        \textbf{Method} & \textbf{Acc} & \textbf{Macro F1} \\
        \midrule
        \multicolumn{3}{l}{\textit{Our methods}} \\
        \rowcolor{blue!8}
        Baseline (InceptionTime) & — & — \\
        MOMENT (linear probe)    & — & — \\
        \rowcolor{green!15}
        \textbf{Ensemble (ours)} & \textbf{—} & \textbf{—} \\
        \midrule
        \multicolumn{3}{l}{\textit{Bake-off SOTA (Ruiz et al., 2021)}} \\
        \textcolor{myred}{MUSE}   & \textcolor{myred}{0.636} & — \\
        ROCKET                   & 0.632 & — \\
        MrSEQL                   & 0.603 & — \\
        \midrule
        \multicolumn{3}{l}{\textit{Recent from-scratch (PRISM 2025)}} \\
        PatchTST                 & 0.417 & — \\
        ROCKET (reimpl.)         & 0.375 & — \\
        PRISM                    & 0.317 & — \\
        \bottomrule
      \end{tabular}
      \end{center}
      \small \textit{— = filled after final training run (see \texttt{train\_best.py})}
    \end{column}
    \begin{column}{0.40\textwidth}
      \textbf{Key findings:}
      \begin{itemize}
        \item Our \textbf{baseline beats all PRISM-2025 methods} — thanks to imbalance handling
        \item MOMENT provides \emph{complementary} predictions
        \item Ensemble target: beat MUSE ($>63.6\%$)
        \item \textbf{No prior work} uses a foundation model on LSST
      \end{itemize}
      \vspace{0.3em}
      \begin{alertblock}{Fair comparison}
        \small Bake-off methods also did not use WeightedSampler.
        Our handling of imbalance accounts for $\sim$8\% gap vs.\ na\"ive training.
      \end{alertblock}
    \end{column}
  \end{columns}
\end{frame}

% ── Slide 10: Per-class results ───────────────────────────────────────────────
\begin{frame}{Per-class Analysis — Ensemble}
  \begin{columns}
    \begin{column}{0.50\textwidth}
      \begin{center}
      \footnotesize
      \begin{tabular}{lrrr}
        \toprule
        Class & Prec & Rec & F1 \\
        \midrule
        \textcolor{mygreen}{16 (SN II)}    & — & — & \textbf{—} \\
        \textcolor{mygreen}{88}            & — & — & \textbf{—} \\
        \textcolor{mygreen}{65}            & — & — & \textbf{—} \\
        \textcolor{mygreen}{92}            & — & — & \textbf{—} \\
        90 (SN Ia, 44\%)                   & — & — & — \\
        42                                 & — & — & — \\
        \textcolor{myred}{64} (24 train)   & — & — & \textbf{—} \\
        \textcolor{myred}{53} (7 train)    & — & — & \textbf{—} \\
        \bottomrule
      \end{tabular}
      \end{center}
      \small \textit{(to be filled after training run)}
    \end{column}
    \begin{column}{0.48\textwidth}
      \textbf{Expected observations:}
      \begin{itemize}
        \item Common classes (16, 88, 65) should score well (F1 $>0.7$)
        \item Class 90 dominant (44\%) but tricky — similar to other SN types
        \item \textbf{Class 64} (24 train samples): likely F1$\approx$0 — few-shot regime
        \item Class 53 (7 samples): precision may be high, recall poor
      \end{itemize}
      \vspace{0.3em}
      \begin{block}{Macro vs Weighted F1}
        \small Macro F1 penalises rare classes equally.\\
        Weighted F1 (by support) is higher and reflects
        real-world relevance better.
      \end{block}
    \end{column}
  \end{columns}
\end{frame}

\section{Conclusion}

% ── Slide 11: Conclusion ─────────────────────────────────────────────────────
\begin{frame}{Conclusion \& Future Work}
  \begin{columns}
    \begin{column}{0.55\textwidth}
      \textbf{Contributions:}
      \begin{enumerate}
        \item \textbf{First application} of a pre-trained time-series
          foundation model (MOMENT-1-large) to LSST classification
        \item \textbf{Competitive results} on LSST:
          Acc=—, Macro F1=— — outperforms all recent
          from-scratch methods; approaches MUSE SOTA (63.6\%)
        \item \textbf{Robust preprocessing} pipeline (NaN handling +
          WeightedRandomSampler) shown to be critical
        \item \textbf{Ensemble} of foundation model +
          complementary architectures via F1-weighted soft-voting
      \end{enumerate}
    \end{column}
    \begin{column}{0.42\textwidth}
      \textbf{Limitations \& Future work:}
      \begin{itemize}
        \item MOMENT phase 3 (full fine-tune) blocked by NaN — needs bf16 training
        \item Classes 64, 53 still poorly classified — data augmentation or few-shot learning
        \item Explore MOMENT-2 (newer version) when available
        \item Test on real LSST survey data (not simulated PLAsTiCC)
      \end{itemize}
      \vspace{0.4em}
      \begin{block}{Take-away}
        \centering
        Foundation model fine-tuning\\
        is viable and competitive\\
        on short multivariate TS.
      \end{block}
    \end{column}
  \end{columns}
\end{frame}

% ── Slide 12: References ─────────────────────────────────────────────────────
\begin{frame}{References}
  \footnotesize
  \begin{thebibliography}{9}
    \bibitem{moment}
      Goswami et al., \textit{MOMENT: A Family of Open Time-series Foundation Models},
      ICML 2024. \href{https://arxiv.org/abs/2402.03885}{arXiv:2402.03885}

    \bibitem{prism}
      Younis et al., \textit{PRISM: Lightweight Multivariate Time-Series Classification},
      2025. \href{https://arxiv.org/abs/2508.04503}{arXiv:2508.04503}

    \bibitem{inception}
      Fawaz et al., \textit{InceptionTime: Finding AlexNet for Time Series Classification},
      DMKD 2020.

    \bibitem{patchtst}
      Nie et al., \textit{A Time Series is Worth 64 Words: Long-term Forecasting with Transformers},
      ICLR 2023.

    \bibitem{rocket}
      Dempster et al., \textit{ROCKET: Exceptionally fast and accurate time series
      classification using random convolutional kernels}, DMKD 2020.

    \bibitem{plasticc}
      The PLAsTiCC Team, \textit{The Photometric LSST Astronomical Time-Series
      Classification Challenge (PLAsTiCC)}, Zenodo 2018.

    \bibitem{bakeoff}
      Ruiz et al., \textit{The Great Multivariate Time Series Classification Bake Off},
      Data Mining and Knowledge Discovery 2021.

    \bibitem{hydra}
      Dempster et al., \textit{Hydra: Competing Convolutional Kernels for Fast and
      Accurate Time Series Classification}, DMKD 2023.
  \end{thebibliography}
\end{frame}

% ── Backup slide ─────────────────────────────────────────────────────────────
\begin{frame}[noframenumbering]{Backup — Training details}
  \begin{center}
  \footnotesize
  \begin{tabular}{llll}
    \toprule
    Model & Epochs & LR & Optimizer \\
    \midrule
    InceptionTime (baseline) & 200 & $10^{-3}$ & AdamW + CosineAnnealing \\
    MOMENT Phase 1           & 50  & $10^{-3}$ & AdamW + CosineAnnealing \\
    MOMENT Phase 2           & 40  & $5\times10^{-4}$ (head) / $5\times10^{-7}$ (enc) & AdamW \\
    PatchTST                 & 150 & $5\times10^{-4}$ & AdamW + CosineAnnealing \\
    InceptionTime-Large      & 200 & $10^{-3}$ & AdamW + CosineAnnealing \\
    FeatureMLP               & 100 & $10^{-3}$ & AdamW + ReduceLROnPlateau \\
    \midrule
    All & \multicolumn{3}{l}{label smoothing=0.1, MixUp $\alpha$=0.3, gradient clip=1.0} \\
    \bottomrule
  \end{tabular}
  \end{center}
\end{frame}

\end{document}
